% Aporte de Juan al informe y al manual del programador: validador de planos.
% Cada \section dice a qué parte del informe o del manual pertenece, para
% poder cortarla y pegarla en el documento integrado del grupo.
% Para compilar solo este archivo: pdflatex validador-y-cargador.tex

\documentclass[11pt,a4paper]{article}
\usepackage[utf8]{inputenc}
\usepackage[T1]{fontenc}
\usepackage[spanish]{babel}
\usepackage[margin=2.5cm]{geometry}
\usepackage{booktabs}
\usepackage{tabularx}
\usepackage{enumitem}
\usepackage{listings}
\usepackage{hyperref}

\lstset{
  basicstyle=\ttfamily\small,
  frame=single,
  breaklines=true,
  columns=fullflexible,
  language=Java
}

\title{Validador de planos y su integraci\'on con el cargador\\
\large Aporte para el informe y el manual del programador}
\author{Juan}
\date{}

\begin{document}
\maketitle
\tableofcontents

% =====================================================================
\part{Para el informe}
% =====================================================================

\section{Dise\~no: decisiones del validador}

El validador (\texttt{ValidadorPlano}, paquete \texttt{construccion}) recibe el
JSON del Anexo B.3 y comprueba las once reglas de B.4. Si el plano es v\'alido
no hace nada; si no, lanza \texttt{IllegalArgumentException} con un mensaje que
empieza por la regla, por ejemplo \texttt{R5: desde la entrada se alcanzan 3 de
39 celdas transitables}. A continuaci\'on, las decisiones que no son obvias
leyendo el c\'odigo.

\subsection*{Decisi\'on 1: rechazar con la regla en el mensaje}

\textbf{Qu\'e elegimos.} \texttt{validar} devuelve \texttt{void} y lanza una
excepci\'on cuyo mensaje empieza por \texttt{Rn:}.

\textbf{Por qu\'e.} El enunciado pide que el rechazo diga cu\'al regla fall\'o, y
los tests necesitan comprobar que se viol\'o \emph{esa} regla y no otra. Con el
prefijo, un test es una l\'inea: \texttt{startsWith("R5:")}.

\textbf{Qu\'e descartamos.} Devolver un \texttt{boolean} (no dice qu\'e regla
fall\'o) y juntar todas las violaciones en una lista (m\'as c\'odigo, y la
especificaci\'on no lo pide). El validador corta en la primera regla que falla.

\subsection*{Decisi\'on 2: un solo BFS para R5 y R9}

\textbf{Qu\'e elegimos.} Un recorrido en anchura con cola desde la entrada,
calculado una sola vez en \texttt{validar}. Deja una matriz de distancias
(\texttt{-1} significa ``no visitada'') que usan dos reglas:
\begin{itemize}[nosep]
  \item \textbf{R5, conexo:} la cantidad de celdas con distancia distinta de
        \texttt{-1} tiene que ser igual a la cantidad de celdas transitables.
  \item \textbf{R9, salida m\'as lejana:} la distancia de la \texttt{S} tiene
        que ser la m\'axima de la matriz (los empates valen).
\end{itemize}

\textbf{Por qu\'e.} R9 necesita distancias m\'inimas, y el BFS las da porque
visita por capas: la primera vez que llega a una celda es por el camino m\'as
corto. Como esa misma matriz responde R5, no hace falta un segundo recorrido.

\textbf{Qu\'e descartamos.} Un DFS para la conexidad: servir\'ia para R5 pero no
para R9, y habr\'ia dos recorridos distintos del mismo plano.

\paragraph{Un detalle de R5.} Las paredes nunca se visitan, as\'i que quedan en
\texttt{-1} incluso en un plano v\'alido. Por eso la regla \emph{cuenta} las
celdas alcanzadas y las compara con las transitables, en vez de buscar si
``queda alg\'un \texttt{-1}'': esa b\'usqueda rechazar\'ia todos los planos.

\subsection*{Decisi\'on 3: porcentajes con enteros}

R6 (al menos 10\,\% de bifurcaciones), R7 (entre 30\,\% y 60\,\% transitables)
y R10 (agua y lava hasta 25\,\%) se comparan multiplicando en vez de dividir,
por ejemplo \texttt{bifurcaciones * 100 < transitables * 10}. As\'i no
aparecen decimales ni errores de redondeo en los bordes: un plano con
exactamente el 10\,\% pasa, como dice la regla.

\subsection*{Decisi\'on 4: el validador no construye nada}

\textbf{Qu\'e elegimos.} El validador solo juzga. El armado del \texttt{Plano}
(terrenos, cofres y monstruos sorteados con los cat\'alogos) lo hace
\texttt{CargadorDePlano}, que llama al validador antes de construir.

\textbf{Por qu\'e.} El validador sirve igual para planos propios y ajenos, y
no necesita conocer los cat\'alogos de cofres y monstruos del grupo. El cargador
recibe los cat\'alogos y el \texttt{Random} por constructor, as\'i se puede
probar con un azar fijo.

\textbf{El costo.} El JSON se lee dos veces, una en el validador y otra en el
cargador. Lo aceptamos a cambio de que cada clase tenga una sola
responsabilidad.

\section{Estructuras de datos}

\begin{table}[h]
\centering
\begin{tabularx}{\textwidth}{lXXX}
\toprule
Estructura & D\'onde se usa & Por qu\'e esa & Costo \\
\midrule
Cola & BFS de \texttt{ValidadorPlano}: conexidad (R5) y salida m\'as lejana (R9)
     & Visita por capas, as\'i la primera llegada a una celda es por el camino
       m\'as corto
     & $O(f \cdot c)$: cada celda entra una vez a la cola y se miran 4 vecinos \\
Matriz \texttt{int[][]} & Distancias desde la entrada
     & Acceso directo por fila y columna; \texttt{-1} marca ``no visitada''
     & $O(1)$ por celda; ocupa $f \cdot c$ enteros \\
\bottomrule
\end{tabularx}
\caption{Estructuras que usa el validador de planos.}
\label{tab:estructuras-validador}
\end{table}

La Tabla~\ref{tab:estructuras-validador} cubre solo lo del validador. La
operaci\'on que m\'as se repite es mirar los cuatro vecinos de una celda, una vez
por celda transitable en el BFS y otra vez en R6.

\section{Pruebas}

\texttt{ValidadorPlanoTest} tiene 15 tests, que se corren con
\texttt{sh gradlew test} (el archivo \texttt{gradlew} no tiene permiso de
ejecuci\'on en este repositorio):

\begin{itemize}[nosep]
  \item 1 test con el plano de ejemplo del Anexo B.3, que debe pasar las once
        reglas (B.5 lo verifica a mano).
  \item 11 tests, uno por regla: cada plano viola \emph{solo} esa regla y el
        test comprueba que el mensaje empieza por ella. Es lo que pide B.7.
  \item 3 tests de JSON inv\'alido: mal formado, que no es un objeto, y sin el
        campo \texttt{plano}.
\end{itemize}

Para R5, R6 y R7 usamos planos chicos armados a mano (5\,$\times$\,7, 5\,$\times$\,5
y 7\,$\times$\,7). Para las dem\'as, el plano de ejemplo con una modificaci\'on
puntual.
% COMPLETAR: captura de gradlew test en verde (Figura X).

\section{Problemas encontrados}

\paragraph{Un test que fallaba por otra regla.} Los primeros tests de R5, R7 y
R10 modificaban el plano de ejemplo y el validador rechazaba por una regla
\emph{distinta} de la esperada. Como las reglas se comprueban en orden
(R1 a R11), una modificaci\'on que rompe R7 puede romper antes R5 o R6. Lo
resolvimos con planos chicos dise\~nados para violar una sola regla, y
ordenando los casos: un plano para R7 tiene que ser conexo y tener
bifurcaciones, porque esas dos reglas se comprueban antes.

\paragraph{Posiciones en base 0 y en base 1.} \texttt{Tablero} y \texttt{Posicion}
usan posiciones de 1 a $n$, igual que el JSON del generador. El validador, en
cambio, convierte a base 0 al leer, porque indexa arreglos de Java. Mientras
el validador solo juzgaba, eso no molestaba; al integrarlo con el cargador
apareci\'o el riesgo de pasar al \texttt{Tablero} una posici\'on corrida una
celda. El cargador evita el problema leyendo \texttt{entrada} y \texttt{salida}
por su cuenta, en base 1.

\paragraph{\texttt{salida} ausente.} El validador acepta un JSON sin la clave
\texttt{salida} cuando \texttt{conSalida} es falso, porque Gson devuelve
\texttt{null} de Java para una clave que no est\'a. El cargador llamaba
\texttt{.isJsonNull()} sobre ese \texttt{null} y tiraba
\texttt{NullPointerException}. Se corrigi\'o con
\texttt{elemento == null || elemento.isJsonNull()} (commit \texttt{3834471}).

\paragraph{El borrador de B.5.} El Anexo B.5 cuenta que un primer borrador del
plano de ejemplo ten\'ia la \texttt{S} a distancia 15 y la celda m\'as lejana a 16,
y que a ojo nadie lo vio. Usamos ese caso para el test de R9.

\section{Limitaciones}

\begin{itemize}
  \item El BFS usa \texttt{java.util.ArrayDeque}. Si el grupo implementa su
        propia \texttt{Cola} en el paquete \texttt{estructuras} (hoy solo tiene
        el \texttt{package-info}), hay que reemplazarla ah\'i. Conviene
        confirmar contra la \'ultima versi\'on del enunciado qu\'e se permite en
        cada paquete.
  \item El validador usa base 0 internamente y el resto del modelo, base 1. Hay
        que unificar o dejarlo documentado.
  \item Todav\'ia no validamos con este validador los tres JSON de
        \texttt{ejemplos/generador} ni un plano de otro grupo (Anexo B.8).
  \item El validador no controla los campos \texttt{version}, \texttt{generador}
        ni \texttt{dificultad} del JSON: B.4 no pide hacerlo.
\end{itemize}

\section{Secci\'on del integrante}

% COMPLETAR por Juan, con sus palabras. Material de base:
Trabajo propio en el validador de planos: las reglas R5 (conexo), R6
(bifurcaciones) y R9 (salida m\'as lejana), con el BFS compartido (commit
\texttt{a962d83}); los tests por regla y los de JSON inv\'alido (mismo commit);
y la correcci\'on del caso de \texttt{salida} ausente en el cargador
(commit \texttt{3834471}).

% =====================================================================
\part{Para el manual del programador}
% =====================================================================

\section{Formato de archivo: el plano}

Es el JSON del Anexo B.3. Lo producen los generadores y lo leen
\texttt{ValidadorPlano} y \texttt{CargadorDePlano}.

\begin{lstlisting}[language={}]
{
  "version": 1,
  "generador": "Grupo 03",
  "parametros": { "version": 1, "dificultad": "MEDIO", "filas": 9,
                  "columnas": 12, "semilla": 7, "cofres": 1,
                  "monstruos": 1, "conSalida": true },
  "entrada": { "fila": 2, "columna": 2 },
  "salida":  { "fila": 8, "columna": 10 },
  "plano": [ "############", "#E..#......#", "...", "############" ]
}
\end{lstlisting}

\begin{table}[h]
\centering
\begin{tabularx}{\textwidth}{lX}
\toprule
Campo & Significado \\
\midrule
\texttt{parametros} & Copia de la entrada del generador. El validador usa
\texttt{filas}, \texttt{columnas}, \texttt{cofres}, \texttt{monstruos} y
\texttt{conSalida}. \\
\texttt{entrada} & D\'onde aparece el h\'eroe. Fila y columna de 1 a $n$. Tiene que
coincidir con la \texttt{E} del plano. \\
\texttt{salida} & D\'onde est\'a la salida, o \texttt{null} si \texttt{conSalida}
es falso. Tiene que coincidir con la \texttt{S}. \\
\texttt{plano} & Una cadena por fila, todas del mismo largo. Es la \'unica fuente
de verdad. \\
\bottomrule
\end{tabularx}
\caption{Campos del JSON de un plano.}
\end{table}

S\'imbolos: \texttt{\#} pared, \texttt{.} piso, \texttt{\~{}} agua, \texttt{\^{}} lava,
\texttt{C} cofre, \texttt{M} monstruo, \texttt{E} entrada, \texttt{S} salida.

\section{Las once reglas}

\begin{table}[h]
\centering
\begin{tabularx}{\textwidth}{lX}
\toprule
Regla & Qu\'e comprueba \\
\midrule
R1 & Dimensiones: \texttt{filas} cadenas de \texttt{columnas} caracteres, ambas $\geq 5$ \\
R2 & Solo los ocho s\'imbolos \\
R3 & Borde de pared \\
R4 & Una \texttt{E}, y una \texttt{S} si \texttt{conSalida}; los campos coinciden \\
R5 & Conexo: desde la \texttt{E} se llega a toda celda transitable \\
R6 & Al menos 10\,\% de las transitables tienen 3 o m\'as vecinos transitables \\
R7 & Entre 30\,\% y 60\,\% de las celdas son transitables \\
R8 & Tantas \texttt{C} y \texttt{M} como \texttt{cofres} y \texttt{monstruos} \\
R9 & La \texttt{S} est\'a a la distancia m\'axima de la \texttt{E} \\
R10 & Agua y lava, como mucho 25\,\% de las transitables \\
R11 & Toda \texttt{M} a distancia Manhattan $\geq 3$ de la \texttt{E} \\
\bottomrule
\end{tabularx}
\caption{Reglas de B.4 que comprueba \texttt{ValidadorPlano}.}
\end{table}

\section{Receta: agregar una regla al validador}

\begin{enumerate}
  \item En \texttt{ValidadorPlano}, escribir un m\'etodo
        \texttt{validarXxx()} que lance
        \texttt{IllegalArgumentException("Rn: ...")}. El mensaje tiene que
        empezar por la regla: los tests lo comprueban con \texttt{startsWith}.
  \item Llamarlo desde \texttt{validar}, \emph{en el lugar correcto}: las reglas
        se comprueban en orden, as\'i que una regla que supone que el plano es
        conexo va despu\'es de R5. Si necesita distancias, recibe la matriz que
        calcula \texttt{calcularDistancias()}.
  \item Agregar en \texttt{ValidadorPlanoTest} un plano que viole \emph{solo}
        esa regla y llamar a \texttt{assertRechaza("Rn", json)}.
  \item Correr \texttt{sh gradlew test}.
\end{enumerate}

\section{Decisiones y trampas}

\begin{itemize}
  \item \textbf{Las posiciones van de 1 a $n$} en \texttt{Posicion},
        \texttt{Tablero}, \texttt{Plano} y en el JSON. El validador usa base 0
        solo para indexar sus arreglos. No le pases al \texttt{Tablero} una
        posici\'on que sali\'o del validador.
  \item \textbf{Las paredes quedan en \texttt{-1}} en la matriz de distancias.
        R5 cuenta las celdas alcanzadas; no busca ``alg\'un \texttt{-1}''.
  \item \textbf{El orden de las reglas importa} para los tests: un plano para
        probar R7 tiene que pasar antes R1 a R6.
  \item \textbf{\texttt{salida} puede faltar o ser \texttt{null}} cuando
        \texttt{conSalida} es falso. Quien lea ese campo tiene que
        contemplar las dos formas: en Gson, una clave ausente devuelve
        \texttt{null} de Java y una clave en \texttt{null} devuelve un
        \texttt{JsonNull}.
  \item \texttt{gradlew} no tiene permiso de ejecuci\'on: usar
        \texttt{sh gradlew test}.
\end{itemize}

\section{Pendientes y bugs conocidos}

\begin{itemize}
  \item Unificar la base de las posiciones del validador con el resto del
        modelo.
  \item Reemplazar \texttt{ArrayDeque} por la \texttt{Cola} propia, si el
        grupo la implementa.
  \item Pasar por el validador los tres JSON de \texttt{ejemplos/generador} y
        el plano de otro grupo, y anotar la regla que falle.
\end{itemize}

\end{document}
